\title{Direct Preference Optimization: Your Language Model is Secretly a Reward Model}

\begin{abstract}
Although large-scale unsupervised language models (LMs) acquire extensive world knowledge and some inferential abilities, exerting fine-grained control over their outputs remains challenging due to the fully unsupervised training regime. Common approaches to improve steerability involve gathering human judgments on the relative quality of model outputs and subsequently fine-tuning the LM to reflect these preferences, typically utilizing reinforcement learning from human feedback (RLHF). However, RLHF entails a complicated and frequently unstable process: it first learns a reward model to capture human preferences and then applies reinforcement learning to adjust the large unsupervised LM to maximize the estimated reward, while constraining divergence from the original model. This work proposes a novel formulation of the reward model in RLHF that admits a closed-form solution for the associated optimal policy. This enables us to address the standard RLHF objective through a straightforward classification loss, bypassing the need for complex optimization procedures. We term our approach \textit{Direct Preference Optimization} (DPO), which is robust, effective, and computationally efficient, obviating both LM sampling during fine-tuning and extensive hyperparameter tuning. Experimental results demonstrate that DPO fine-tunes LMs to human preferences on par with, or superior to, prevalent techniques. Remarkably, DPO achieves better sentiment control than PPO-based RLHF and attains equivalent or higher performance in summarization and single-turn dialogue tasks, all while offering a much simpler and more efficient implementation.
\end{abstract}

\section{Introduction}
Large-scale unsupervised language models (LMs), when trained on expansive datasets, often exhibit remarkable and unexpected abilities~\citep{chowdhery2022palm, brown2020language, touvron2023llama, bubeck2023sparks}. However, the underlying training data is produced by humans with diverse intentions, priorities, and varying expertise. Imitating all such behaviors is not always desirable; for instance, we may require an AI coding assistant to \textit{recognize} common programming errors for the purpose of rectifying them, but during code generation, our preference is for the model to reflect the (possibly rare) superior coding practices found in the dataset. Likewise, a language model should ideally be \textit{cognizant} of widespread misconceptions—such as one believed by 50\% of individuals—yet we do not wish for the model to perpetuate that falsehood in half the instances it is queried about it. Therefore, isolating and eliciting the \emph{intended outputs and behaviors} from the broader range of the model’s \textit{knowledge and competencies} is vital for developing AI systems that are reliable, effective, and controllable~\citep{ouyang2022training}. Traditional strategies frequently use reinforcement learning (RL) to align language models with human preferences, but as we will demonstrate, the RL-based loss optimized in these frameworks can be equivalently expressed—and directly optimized—through a straightforward binary cross-entropy objective, thereby streamlining the entire preference learning process.

In broad terms, contemporary approaches impart the required behaviors into a language model by relying on curated datasets composed of human preferences, thereby representing behaviors regarded as safe and useful by people. This phase of learning human preferences follows an initial stage where the model undergoes extensive unsupervised pre-training on large corpora of text. While supervised fine-tuning on examples of preferred human behavior is a direct approach to learning from preferences, the leading methods utilize reinforcement learning from human (or AI-generated) feedback (RLHF/RLAIF; \citep{christiano2017deep,bai2022constitutional}). RLHF methodologies involve first fitting a reward model to human preference data, followed by applying RL to adjust a language model’s policy so that it produces high-reward responses without straying too far from the original model’s behavior. Despite RLHF’s success in producing models excelling at both conversational and coding tasks, its training regime is significantly more elaborate than that of supervised approaches, requiring multiple language models and integrating policy sampling within the training loop, which leads to substantial increases in computational overhead.

This work introduces a method to directly align language models with human preferences, eliminating the need for explicit reward modeling or reinforcement learning. We present \textit{Direct Preference Optimization} (DPO), a technique that implicitly optimizes for the same reward-maximization objective as standard RLHF methods—incorporating a KL-divergence constraint—while maintaining ease of implementation and training. Conceptually, DPO updates serve to boost the ratio of the log probabilities assigned to favored responses over less-preferred ones, integrating a per-example, adaptive importance weight that mitigates the degradation in model quality observed when simply applying a naïve probability ratio objective. Similar to existing algorithms, DPO leverages a theoretical framework for preference modeling, such as the Bradley-Terry model \cite{bradley1952rankanalysis}, which evaluates the correspondence between a candidate reward function and empirical preference annotations. Distinctively, while other methods use this framework to define a loss for fitting a reward model, then subsequently train a policy to optimize this learned reward, DPO re-parameterizes the preference loss directly in terms of the policy. With access to a dataset of human-labeled preferences, DPO enables direct optimization of the policy via a standard binary cross-entropy loss, yielding an optimal policy with respect to a latent reward function inferred from the data.

Our primary contribution is Direct Preference Optimization (DPO), a straightforward algorithm that dispenses with reinforcement learning for preference-driven language model training. Empirical results demonstrate that DPO performs on par with leading alternatives, such as PPO-based RLHF, on a variety of preference-based tasks including sentiment control, summarization, and dialogue, evaluated on language models scaling up to 6B parameters.

Self-supervised language models, when scaled up, are capable of addressing some tasks in a zero-shot manner \citep{radford2019language} or with minimal prompting in few-shot settings \citep{gpt3,megatron,chowdhery2022palm}. Nevertheless, their downstream effectiveness and alignment with user intentions are greatly enhanced when further fine-tuned on corpora containing explicit instructions and human-crafted completions \citep{mishra-etal-2022-cross,sanh2022multitask,chung2022scaling,thoppilan2022lamda}. This technique, known as `instruction-tuning,' allows LLMs to better generalize to instructions that were not present in their training datasets and generally improves their practicality \citep{chung2022scaling}. Although instruction tuning has demonstrated considerable success, collecting \textit{relative} assessments of output quality from humans is frequently less demanding than sourcing expert-labeled demonstrations. As a result, recent efforts have involved refining LLMs using datasets constructed from human preference judgments, leading to advances in capabilities such as translation \citep{kreutzer-etal-2018-reliability}, summarization \citep{stiennon2022learning,ziegler2020finetuning}, story generation \citep{ziegler2020finetuning}, and executing instructions \citep{ouyang2022training,ramamurthy2023is}. Typically, these methods proceed by first calibrating a neural reward function to fit observed human preferences, often modeled via frameworks like the Bradley-Terry model \citep{bradley1952rankanalysis}, and subsequently adjusting the language model’s parameters to maximize this reward through reinforcement learning strategies such as REINFORCE \citep{williams1992reinforce}, proximal policy optimization (PPO; \cite{schulman2017proximal}), or their derivatives \citep{ramamurthy2023is}. In parallel, another research stream utilizes instruction-following LLMs trained with human feedback to synthesize additional preference data targeting particular characteristics, such as safety or non-maliciousness \citep{bai2022constitutional}, relying on loosely structured human oversight via textual rubrics for annotating model outputs. This intersection merges two key research directions: one focusing on the use of reinforcement learning for tuning language models towards diverse goals~\citep{Ranzato2015SequenceLT,paulus2018a,wu2018learning}, and another centered on frameworks for learning from human preference signals \citep{christiano2017deep,kupcsik2018learning}. Although harnessing relative human judgments is advantageous, applying reinforcement learning at scale for LLM fine-tuning remains a significant logistical barrier; the present work introduces a theoretically principled alternative for optimizing over relative preferences that circumvents the use of RL.

Beyond linguistic applications, the study of policy learning from preferences has attracted attention in both bandit and reinforcement learning frameworks, resulting in various proposed methodologies. When action preferences or rankings, instead of explicit rewards, are utilized for learning in contextual bandit problems, the setup is referred to as contextual dueling bandit (CDB; \cite{yue2012karmed,dudik2015contextual}). Since absolute reward signals are absent in this scenario, theoretical investigations into CDBs replace the concept of an optimal policy with that of a \textit{von Neumann winner}, defined as a policy that achieves an expected win rate of no less than 50\% against any other policy \citep{dudik2015contextual}. Notably, CDB frameworks generate preference labels interactively in an online manner, whereas learning from human preferences often relies on a pre-collected, finite dataset comprising offline preference-annotated action pairs \citep{yan2022human}. Likewise, \textit{preference-based RL} (PbRL) draws upon binary preferences derived from an unknown scoring function in lieu of explicit rewards \citep{BusaFekete2014,ruiz2023dueling}. A range of PbRL algorithms have been introduced, with some capable of leveraging off-policy preference datasets. These methods typically consist of a two-stage process: first, the estimation of the underlying scoring (reward) function, followed by policy optimization based on the estimated model \citep{jain2013learning,BusaFekete2014,christiano2017deep,sadigh2017active,kupcsik2018learning}. In contrast, we introduce a unified approach that trains the policy in a single stage to directly adhere to observed preferences.

\section{Preliminaries}\label{section:prelims}

Here, we provide an overview of the RLHF framework as described by \citeauthor{ziegler2020finetuning}, and subsequently employed in \citep{stiennon2022learning, bai2022training, ouyang2022training}. The procedure is generally organized into three stages: 1) supervised fine-tuning (SFT), 2) collection of preference data and reward modeling, and 3) optimization with reinforcement learning.

\textbf{SFT}: The RLHF process customarily starts with adapting a pre-trained language model through supervised learning on a curated dataset that is relevant to the target applications (such as dialogue or summarization tasks), yielding the model $\pi^\text{SFT}$. 

\textbf{Reward Modeling Phase}: During the next step, prompts $x$ are fed to the SFT model to generate pairs of responses $(y_1, y_2)\sim \pi^\text{SFT}(y \mid x)$. Human annotators then evaluate these response pairs and express their preferences, represented as $y_w\succ y_l \mid x$—here, $y_w$ is the favored and $y_l$ the less favored answer among $(y_1, y_2)$. These expressed preferences are considered to be governed by a latent reward function $r^*(y, x)$ that is inaccessible to us. Several statistical models have been proposed to capture this preference mechanism, with the Bradley-Terry (BT) model \cite{bradley1952rankanalysis} being a widely adopted one (though the Plackett-Luce family of ranking models \citep{plackett1975analysis, luce2012individual} is also compatible if multiple ranked completions are available). The BT model postulates that the distribution of human preferences $p^*$ may be represented as follows:

Given access to a fixed dataset of pairwise comparisons $\mathcal{D}=\bigl\{x^{(i)}, y_w^{(i)}, y_l^{(i)}\bigr\}_{i=1}^N$, which are assumed to be sampled according to $p^*$, one may parameterize a reward function $r_{\phi}(x, y)$. The parameters of this model can be fitted using maximum likelihood estimation. When formulated as a binary classification task, the associated negative log-likelihood loss can be expressed as:

\begin{equation}\label{eq:reward_model}
    \mathcal{L}_R(r_{\phi}, \mathcal{D}) = -\mathbb{E}_{(x, y_w, y_l)\sim \mathcal{D}}\bigl[\log \sigma(r_{\phi}(x, y_w)- r_{\phi}(x, y_l))\bigr]
\end{equation}

where $\sigma$ denotes the logistic sigmoid function. In large language models (LMs), $r_{\phi}(x, y)$ is commonly initialized from the SFT model $\pi^\text{SFT}(y \mid x)$, augmented by a linear projection atop the final transformer representation to yield a scalar reward prediction \cite{ziegler2020finetuning}. To mitigate variance in the learned reward, earlier works often apply normalization so that  $\mathbb{E}_{x,y\sim \mathcal{D}}\left[r_\phi(x, y)\right] = 0$ for each $x$.

\textbf{RL Fine-Tuning Phase}: In this stage, the reward model is leveraged to score outputs produced by the language model. As outlined by previous studies~\citep{jaques2017sequence, jaques2020human}, the optimization procedure is formulated as:

\begin{equation}\label{eq:RL}
\max_{\pi_{\theta}}  \mathbb{E}_{x\sim \mathcal{D}, y\sim \pi_{\theta}(y \mid x)}\bigl[r_{\phi}(x, y)\bigr] - \beta\mathbb{D}_{\textrm{KL}}\bigl[\pi_{\theta}(y\mid x)\mid \mid \pi_\text{ref}(y\mid x)\bigr],
\end{equation}

where $\beta$ determines the strength of regularization toward the reference policy $\pi_\text{ref}$, which is typically chosen as the SFT model $\pi^\text{SFT}$. The policy $\pi_\theta$ is typically initialized from this same SFT model. Incorporating this regularization is critical—it discourages the model from straying excessively far from the distribution where the reward is reliable, helps preserve generative diversity, and reduces the risk of mode collapse on singular high-reward completions. Since text generation is discrete, direct differentiation is not feasible, so policy optimization is carried out using reinforcement learning algorithms. Conventionally \citep{ziegler2020finetuning, stiennon2022learning, bai2022training, ouyang2022training}, the reward is computed as ${r(x, y) = r_{\phi}(x, y) -\beta (\log \pi_{\theta}(y\mid x) - \log \pi_\text{ref}(y\mid x))}$ and the PPO algorithm \cite{schulman2017proximal} is applied to maximize the objective.

\section{Direct Preference Optimization}\label{sec:DPO}

With the practical difficulties of applying standard reinforcement learning algorithms to large-scale models (such as LMs), we aim to develop a streamlined, preference-driven approach to policy learning. Unlike previous RLHF frameworks, which involve an explicit reward model followed by a reinforcement learning step, our strategy is to use a specific reward model structure that allows the direct, closed-form extraction of its optimal policy, obviating the need for RL loops. As we elaborate below, our principal idea is to utilize an explicit, analytical correspondence between reward models and their optimal policies, effectively reparameterizing the loss over reward functions into a loss directly over policies. Through this change-of-variables, the need to fit a separate reward model is eliminated, while continuing to respect established models of human preferences like the Bradley-Terry model. Ultimately, the policy network serves a dual function—encoding both the language model's behavior and (implicitly) its reward evaluation.

\textbf{Derivation of the DPO objective.} Our starting point is the canonical RL objective expressed in Eq.~\ref{eq:RL}, which involves an arbitrary reward function $r$. Following established approaches~\citep{peters2007reinforcement, peng2019advantage, korbak2022reinforcement, go2023aligning}, it can be directly established that the solution to the KL-regularized reward maximization problem (see Eq.~\ref{eq:RL}) assumes the following form:

\begin{equation}\label{eq:op_policy}
    \pi_r(y\mid x) = \frac{1}{Z(x)}\pi_\text{ref}(y\mid x)\exp\left(\frac{1}{\beta}r(x, y)\right),
\end{equation}

where the normalization factor $Z(x) =\sum_{y}\pi_\text{ref}(y\mid x)\exp\left(\frac{1}{\beta}r(x, y)\right)$ denotes the partition function. For a full exposition of the derivation, refer to Appendix~\ref{app:derivation1}. However, even when employing the MLE estimator $r_{\phi}$ for the authentic reward function $r^*$, computing the partition function $Z(x)$ remains computationally demanding~\citep{korbak2022reinforcement, go2023aligning}, which complicates the direct application of this policy form. Nonetheless, we may invert Eq.~\ref{eq:op_policy} so that the reward is instead written in terms of the optimal policy $\pi_r$, the reference policy $\pi_\text{ref}$, and the partition function $Z(\cdot)$. Concretely, by applying the logarithm to both sides of Eq.~\ref{eq:op_policy} and performing algebraic manipulations, we find:

\begin{equation}\label{eq:main_eq}
    r(x,y) =\beta \log \frac{\pi_r(y\mid x)}{\pi_\text{ref}(y\mid x)} + \beta \log Z(x).
\end{equation}

This reparametrization can be applied to both the ground-truth reward $r^*$ and its corresponding optimal policy $\pi^*$. The Bradley-Terry model conveniently utilizes only reward differences between output candidates: ${p^*(y_1 \succ y_2 \mid x) = \sigma(r^*(x, y_1) - r^*(x, y_2))}$. By substituting the form from Eq.~\ref{eq:main_eq} into the Bradley-Terry preference model (Eq.~\ref{eq:bradley-terry}), the dependence on the partition function $Z(x)$ disappears, allowing the human preference likelihood to be stated exclusively via the optimal policy $\pi^*$ and reference policy $\pi_\text{ref}$. Therefore, under the Bradley-Terry model, the optimal RLHF policy $\pi^*$ obeys the preference equation:

\begin{equation}\label{eq:objective}
    p^*(y_1\succ y_2 \mid x)=\frac{1}{1 + \exp\left(\beta \log \frac{\pi^*(y_2\mid x)}{\pi_\text{ref}(y_2\mid x)} - \beta \log \frac{\pi^*(y_1\mid x)}{\pi_\text{ref}(y_1\mid x)}\right)}
\end{equation}

A complete derivation can be found in Appendix~\ref{app:derivation2}. While the above makes use of the Bradley-Terry setup, analogous forms can be derived under more general ranking frameworks, such as the Plackett-Luce model~\citep{plackett1975analysis, luce2012individual}, discussed further in Appendix~\ref{app:plackett_luce_models}.

Since we have now represented the likelihood of human preference data directly in terms of the optimal policy, we may pose a maximum likelihood estimation task for a parameterized policy $\pi_\theta$. Mirroring the reward-based modeling strategy (e.g., Eq.~\ref{eq:reward_model}), this leads to the following policy learning objective:

\begin{equation}\label{eq:optimum_model}
    \mathcal{L}_\text{DPO}(\pi_{\theta}; \pi_\text{ref}) = -\mathbb{E}_{(x, y_w, y_l)\sim \mathcal{D}}\left[\log \sigma \left(\beta \log \frac{\pi_{\theta}(y_w\mid x)}{\pi_\text{ref}(y

In this approach, we represent the implicit reward through an alternative parameterization, for which the optimal policy corresponds directly to $\pi_\theta$. Additionally, since the method amounts to fitting a reparameterized Bradley-Terry model, it inherits certain theoretical guarantees—such as consistency, provided that the preference data distribution meets appropriate conditions \cite{bong2022generalized}. Section~\ref{sec:theory} elaborates further on the theoretical aspects of DPO and compares them to related approaches.

\textbf{What effect does the DPO update have?} To build intuition about the mechanics of DPO, we examine the gradient of the objective function $\mathcal{L}_\text{DPO}$. The gradient with respect to the model parameters $\theta$ is given by:

\begin{multline*}\label{eq:gradient}
    \nabla_\theta \mathcal{L}_\text{DPO}(\pi_\theta;\pi_\text{ref}) = \\ -\beta\mathbb{E}_{(x, y_w, y_l) \sim \mathcal{D}} \bigg[\underbrace{\sigma(\hat{r}_\theta(x, y_l) - \hat{r}_\theta (x, y_w))}_\text{greater influence when reward assignment is incorrect}\bigg[\underbrace{\nabla_\theta\log \pi(y_w \mid x)}_\text{encourage $y_w$} - \underbrace{\nabla_\theta\log\pi(y_l \mid x)}_\text{discourage $y_l$}\bigg]\bigg],
\end{multline*}

Here, $\hat{r}_\theta(x, y) = \beta \log \frac{\pi_\theta(y \mid x)}{\pi_\text{ref}(y \mid x)}$ defines the implicit reward associated with the model $\pi_\theta$ relative to the reference $\pi_\text{ref}$ (see Section~\ref{sec:theory} for further detail). Conceptually, the gradient of $\mathcal{L}_\text{DPO}$ serves to raise the probability of the favored responses $y_w$ while lowering that of the less preferred responses $y_l$. Crucially, the magnitude of each example's contribution is modulated by the extent to which the implicit reward function $\hat{r}_\theta$ incorrectly favors the dispreferred completion, scaled by $\beta$—thereby reflecting both model misranking and the impact of the KL constraint. Empirically, we find this dynamic weighting to be vital; omitting the weighting factor in a naïve variant leads to failure modes in the language model (see Appendix Table~\ref{tab:unlikelihood_generations}).

\textbf{DPO overview.} The typical DPO workflow proceeds as follows: (1) For each prompt $x$, generate candidate completions $y_1, y_2$ sampled from $\pi_\text{ref}(\cdot \mid x)$, then use human annotations to create the preference dataset $\mathcal{D} = \{x^{(i)}, y_w^{(i)}, y_l^{(i)}\}_{i=1}^N$; and (2) update the language model $\pi_\theta$ by minimizing $\mathcal{L}_\text{DPO}$, leveraging the fixed $\pi_\text{ref}$, the preference data $\mathcal{D}$, and the specified $\beta$ parameter. In real-world scenarios, it is generally preferable to utilize publicly available preference datasets, instead of producing completions and collecting preference annotations anew. If such datasets were collected using $\pi^\text{SFT}$, we set $\pi_\text{ref} = \pi^\text{SFT}$ by default. Alternatively, if $\pi^\text{SFT}$ is not accessible, we determine $\pi_\text{ref}$ by fitting it to maximize the likelihood of the preferred completions $(x, y_w)$ observed in the data, i.e., ${\pi_\text{ref} = \argmax_{\pi}\mathbb{E}_{x, y_w \sim \mathcal{D}}\left[\log \pi(y_w \mid x)\right]}$. This approach helps address the distributional gap between the genuine (but unobserved) reference distribution and the practical choice of $\pi_\text{ref}$ within DPO. For a comprehensive description of implementation specifics and hyperparameter settings, refer to Appendix~\ref{app:implementation}.

\section{Theoretical Analysis of DPO}
Within this section, we deepen our interpretation of the DPO algorithm, present theoretical insights supporting the method, and connect the strengths of DPO to shortcomings inherent in actor-critic frameworks for RLHF, such as PPO~\cite{schulman2017proximal}.

\label{sec:theory}

\subsection{Your Language Model Is Secretly a Reward Model}
DPO unifies both implicit reward modeling and policy optimization into a single maximum likelihood training step, without needing to explicitly learn a reward function or engage in separate reinforcement learning procedures. Crucially, the loss function in Eq.~\ref{eq:main_eq} aligns with a Bradley-Terry formulation, where rewards are given by $r^*(x, y) = \beta \log\frac{\pi^*_\theta(y \mid x)}{\pi_\text{ref}(y \mid x)}$. By optimizing the model $\pi_\theta$ in this way, we effectively mirror the reward model maximization problem from Eq.~\ref{eq:reward_model}, provided we apply the corresponding change of variables. Here, we systematically establish the theoretical foundation for this reparameterization, show that it does not limit the expressive capacity of learned reward functions, and demonstrate that it permits perfect recovery of the optimal policy. To that end, we start by introducing an equivalence relation over reward functions.

\begin{definition}
Two reward functions $r(x, y)$ and $r'(x, y)$ are defined as equivalent if and only if there exists some function $f$ such that ${r(x, y)-r'(x, y) = f(x)}$.     
\end{definition}

One can readily verify that this definition indeed characterizes an equivalence relation, thereby dividing the set of all reward functions into distinct equivalence classes. The following lemmas hold:

\begin{lemma}\label{lemma:same_prefrence}
Within the Plackett-Luce model, as well as the Bradley-Terry preference framework, reward functions that belong to the same equivalence class generate identical preference distributions.
\end{lemma}

\begin{lemma}\label{lemma:same_policy}
Any two reward functions drawn from a common equivalence class will yield the same optimal policy under the corresponding constrained RL setting.
\end{lemma}

We postpone the (elementary) proofs to Appendix \ref{app:lemma1}. The first lemma is an instance of a well-recognized under-specification phenomenon inherent in the Plackett-Luce family \cite{plackett1975analysis}. This ambiguity necessitates the use of additional identifiability conditions in order to make any rigorous statements about maximum likelihood estimates arising from Eq. \ref{eq:reward_model} \cite{bong2022generalized}. The second lemma clarifies that all reward functions within a particular equivalence class prescribe the same optimal policy; thus, our downstream objective can focus on identifying any representative from the target class. In Appendix~\ref{app:thm1}, we establish the following theorem:

\begin{theorem}\label{thm:main}
Assuming mild technical conditions, every equivalence class of reward functions compatible with the Plackett-Luce (and specifically Bradley-Terry) models can be characterized via the reparameterization ${r(x, y) = \beta \log \frac{\pi(y\mid x)}{\pi_\text{ref}(y\mid x)}}$, for a suitable model $\pi(y\mid x)$ and a fixed reference policy $\pi_\text{ref}(y \mid x)$.
\end{theorem}

\begin{sproof}
Take an arbitrary reward function $r(x, y)$, which determines its corresponding optimal policy $\pi_r(y \mid x)$, as defined in Eq. \ref{eq:op_policy}. We now demonstrate that a reward function within the same equivalence class as $r$ admits a representation using the stated reparameterization. Define the projection mapping $f$ as
\begin{equation}
    f(r; \pi_\text{ref}, \beta)(x, y) = r(x, y) - \beta\log\sum_{y}\pi_\text{ref}(y\mid x)\exp\left(\frac{1}{\beta}r(x, y)\right)
\end{equation}
Here, the operator $f$ effectively normalizes the reward function by subtracting the log-partition function associated with $\pi_r$. Because the normalization term depends solely on $x$, $f(r; \pi_\text{ref}, \beta)(x, y)$ and $r(x, y)$ belong to the same equivalence class. Lastly, substituting the expression from Eq.~\ref{eq:main_eq} (which universally holds for any reward function), we get $f(r; \pi_\text{ref}, \beta)(x, y) = \beta \log \frac{\pi_r(y\mid x)}{\pi_\text{ref}(y\mid x)}$. Thus, this projection yields a canonical representative in the equivalence class of $r$, and this reparameterization involves no loss of expressiveness for the reward model.
\end{sproof}

Theorem~\ref{thm:main} can also be interpreted as precisely determining which reward function is chosen within each equivalence class by the DPO reparameterization—specifically, the reward function that fulfills:

\begin{equation}\label{eq:lag_p}
     \sum_{y}\underbrace{\pi_\text{ref}(y\mid x)\exp\left(\frac{1}{\beta}r(x, y)\right)}_{=\pi(y\mid x)\text{, using Thm.~\ref{thm:main} reparam.}} = 1,
\end{equation}

This condition ensures that $\pi(y\mid x)$ constitutes a legitimate probability distribution, where all probabilities are positive and total to one. Further, in light of Eq.~\ref{eq:op_policy}, it becomes evident that Eq.~\ref{eq:lag_p} serves as the partition function associated with the optimal policy derived from the reward function $r(x, y)$. The principal insight underlying the DPO algorithm is that, by imposing specific constraints on the otherwise under-determined Plackett-Luce (and in particular, Bradley-Terry) family of preference-based models, we can retain the expressive power of the set of admissible reward functions, while also ensuring that the optimal policy in Eq. \ref{eq:op_policy} becomes explicitly solvable for every prompt $x$.

\subsection{Instability of Actor-Critic Algorithms} Our framework also facilitates the identification of instability issues inherent in typical actor-critic algorithms applied in RLHF, such as PPO. We consider the standard RLHF setup, focusing on the RL fine-tuning procedure as discussed in Section \ref{section:prelims}. There is a natural correspondence with the control as inference framework \cite{levine2018reinforcement} for the constrained RL problem described by \ref{eq:RL}. Assuming a model $\pi_{\theta}(y\mid x)$ parameterized by $\theta$, we minimize the KL divergence $\mathbb{D}_{\text{KL}}[\pi_{\theta}(y|x) \mid \mid \pi^*(y\mid x)]$, where $\pi^*$ is the optimal policy induced by the reward model $r_{\phi}(y, x)$, as defined in Eq. \ref{eq:optimum_model}. After some manipulation, we obtain the following optimization objective:

\begin{equation}\label{eq:AC}
    \max_{\pi_{\theta}}\mathbb{E}_{\pi_{\theta}(y\mid x)}\bigg[\underbrace{r_{\phi}(x, y) -\beta\log\sum_{y}\pi_\text{ref}(y\mid x)\exp\left(\frac{1}{\beta}r_{\phi}(x, y)\right)}_{f(r_{\phi}, \pi_\text{ref}, \beta)} - \underbrace{\beta\log\frac{\pi_{\theta}(y\mid x)}{\pi_\text{ref}(y\mid x)}}_{\text{KL}}\bigg]
\end{equation}

The objective here mirrors that optimized in earlier studies 
\citep{ziegler2020finetuning, stiennon2022learning, bai2022training, ouyang2022training}, utilizing a reward from the $r_{\phi}$ class that is equivalent to that used in DPO. Within this framework, the normalization component of $f(r_{\phi}, \pi_\text{ref}, \beta)$ can be viewed as the soft value function corresponding to the reference policy $\pi_\text{ref}$. Although this normalization factor does not influence the final optimum, its absence could introduce substantial variance into the policy gradient, potentially destabilizing the learning process. One way to address the normalization term is by employing a learned value function, though optimizing this function itself can be challenging. Another approach previously adopted has been to standardize rewards using a human completion as a baseline, which acts as a Monte-Carlo estimate (with a single sample) for the normalization factor. By contrast, the DPO reparameterization constructs a reward function that eliminates the requirement for any such baseline.

\section{Experiments}
Here, we provide an empirical assessment of DPO's performance in learning policies directly from preferences. Our initial analysis uses a carefully controlled text generation scenario to examine the efficiency with which DPO balances reward maximization against the need to minimize KL-divergence from the reference policy, in comparison to widely-used preference-based learning techniques such as PPO. We then extend our evaluation to consider DPO on larger model scales and more complex RLHF problems, including tasks in summarization and conversational agents. Remarkably, we observe that DPO, with minimal hyperparameter adjustment, typically achieves results on par with or superior to competitive methods like RLHF using PPO, as well as surpassing strategies that select the best trajectory from $N$ samples based on a trained reward function. Prior to discussing these empirical findings, we outline our experimental methodology; further specifics can be found in Appendix~\ref{app:exp_details}.

\textbf{Tasks.} Our study examines three open-domain text generation tasks. For all tasks, learning algorithms derive a policy based on a dataset of pairwise preferences $\mathcal{D}=\bigl\{x^{(i)}, y_w^{(i)}, y_l^{(i)}\bigr\}_{i=1}^N$. In \textbf{controlled sentiment generation}, the input $x$ consists of an initial fragment of a movie review taken from the IMDb dataset \cite{maas-EtAl:2011:ACL-HLT2011}, and the model is required to generate a continuation $y$ that expresses positive sentiment. To ensure evaluation objectivity, preference pairs for this task are \textit{automatically generated} over model outputs by employing a pre-trained sentiment classifier, such that $p(\text{positive}\mid x,y_w)>p(\text{positive}\mid x,y_l)$. For supervised fine-tuning (SFT), GPT-2-large is trained to convergence on the IMDb train set reviews; additional details can be found in App~\ref{app:sentiment_details}. In the \textbf{summarization} task, $x$ represents a Reddit forum post, and the goal is for the policy to output a summary $y$ that encapsulates the main points. Consistent with prior research, we utilize the Reddit TL;DR summarization corpus \citep{volske-etal-2017-tl} in conjunction with human preferences curated by \citeauthor{stiennon2022learning}. The SFT baseline is a model further trained on human-crafted Reddit summaries,\footnote{\url{https://huggingface.co/CarperAI/openai_summarize_tldr_sft}} using the TRLX \citep{leandro_von_werra_2023_7790115} RLHF toolkit. Note that the human preference data originates from outputs of a distinct, albeit similarly fine-tuned, SFT model collected by \citeauthor{stiennon2022learning}. Lastly, the \textbf{single-turn dialogue} task frames $x$ as a user query—these range from technical questions about astrophysics to personal advice requests. The objective is for the model to generate a relevant and engaging response $y$. Here, we employ the Anthropic Helpful and Harmless dialogue dataset \citep{bai2022training}, which features 170k interactions between humans and AI assistants; each interaction ends with two model-generated responses and a label indicating the human’s preference. Since there is no dedicated pre-trained SFT model for this dataset, we construct one by fine-tuning an existing language model exclusively on preferred response completions.

\textbf{Evaluation.} We employ two complementary strategies for evaluation in our experiments. To scrutinize the ability of each algorithm to optimize the constrained reward maximization objective, we assess performance in the controlled sentiment generation scenario by plotting each algorithm’s trade-off curve between the obtained reward and its KL-divergence from a reference policy. Since the true reward function (implemented as a sentiment classifier) is available in this setting, the frontier can be directly calculated. In contrast, the true reward function is not observable in practical applications. Therefore, for summarization and single-turn dialogue tasks, we rely on \textit{win rate} comparisons versus a baseline policy, using GPT-4’s judgments as a proxy for human evaluations of both summary quality and response usefulness. For summarization, baseline comparisons use the reference summaries from the test split; for dialogue, the baseline corresponds to the annotator-preferred answer in the test data. Although prior research has indicated that large language models can outperform standard automated metrics in evaluation tasks \citep{Chen2023ExploringTU}, we validate our dependence on GPT-4 by conducting a human evaluation study, presented in Sec.~\ref{sec:human-judgments}. Our results indicate that GPT-4’s ratings are highly consistent with human judgments, and that the agreement rate between humans and GPT-4 often matches or surpasses agreement among different human annotators.

\textbf{Methods.} Beyond DPO, we benchmark a diverse set of established methodologies for aligning language model behavior with human preference data. For the summarization domain, we apply zero-shot prompting to \textbf{GPT-J} \citep{gpt-j}, while for single-turn dialogue we adopt a 2-shot prompting setup using \textbf{Pythia-2.8B} \citep{biderman2023pythia}. We further evaluate a standard \textbf{SFT} model and a variant termed \textbf{Preferred-FT}, which is fine-tuned in a supervised fashion on the preferred completion $y_w$, produced by either the SFT model (for sentiment control and summarization tasks) or a generic LM (in dialogue experiments). Another supervised-style approach we include is \textbf{Unlikelihood}~\citep{welleck2019neural}, which explicitly trains the model to increase the likelihood of $y_w$ while concurrently decreasing the likelihood assigned to $y_l$; an optional coefficient $\alpha\in[0,1]$ modulates the strength of the unlikelihood loss term. Additionally, we incorporate \textbf{PPO} \citep{schulman2017proximal}, employing a reward model trained on preference data, as well as \textbf{PPO-GT}, an oracle variant leveraging the ground-truth reward in the sentiment control setup. Our sentiment control experiments use two PPO-GT versions: a standard implementation \cite{leandro_von_werra_2023_7790115} and an adjusted one with reward normalization and enhanced hyperparameter tuning (the same tuning is applied when running PPO with learned rewards). Lastly, we introduce the \textbf{Best of $N$} baseline, which samples $N$ outputs from the SFT model (or Preferred-FT for dialogue) and selects the output with the highest score per a learned reward function from preference data. This approach, though highly effective, is computationally demanding for even moderate $N$, since it requires $N$ samples at test time for every input.

\subsection{How well can DPO optimize the RLHF objective?}

The KL-constrained reward maximization objective, central to conventional RLHF approaches, is designed to encourage higher rewards without allowing the policy to drift excessively from the reference policy. Thus, when evaluating different methods, it is essential to jointly consider both the magnitude of reward obtained and the resulting KL divergence; pursuing greater reward at the expense of substantially increased KL is often not advantageous. Figure~\ref{fig:frontier-tldr-main} presents the trade-off curve between reward and KL divergence across various algorithms within the sentiment context. Each algorithm undergoes several independent training sessions, each session varying the policy's conservativeness hyperparameter (target KL $\in\{3,6,9,12\}$ for PPO, $\beta \in \{0.05,0.1,1,5\}$, and $\alpha\in\{0.05,0.1,0.5,1\}$ for unlikelihood, while different random seeds are used for preferred-FT). Altogether, this produces a total of 22 distinct training runs. Following every 100 steps up to convergence, we assess the resulting policies against a held-out set of prompts, calculating both the mean reward according to the ground truth reward function and the mean sequence-level KL\footnote{That is, the sum of the per-timestep KL-divergences.} relative to the reference policy, i.e., $\text{KL}\left(\pi\mid \mid \pi_\text{ref}\right)$. The results indicate that DPO establishes a clearly superior frontier, yielding the maximal achievable reward while also maintaining minimal KL. This outcome is particularly significant for several reasons. Notably, both DPO and PPO target the same optimization criterion, yet DPO attains substantially greater efficiency; DPO achieves a reward/KL balance that universally surpasses that of PPO. Furthermore, DPO outperforms PPO on this frontier, \emph{even in scenarios where PPO is provided access to the true reward function} (PPO-GT).

\subsection{Can DPO scale to real preference datasets?}
\label{sec:dpo-real-datasets}
We proceed to assess the effectiveness of DPO fine-tuning when applied to tasks involving summarization and single-turn dialogue. In the context of summarization, it is well established that automated evaluation metrics like ROUGE may not align closely with human judgment~\citep{stiennon2022learning}, and previous research has demonstrated that tuning language models using PPO based on human preference data can yield more preferable summaries. For our analysis, we generate completions using various approaches on the test portion of the TL;DR summarization dataset and calculate the average rate at which these outputs are preferred over reference completions within the test set. For every method, completions are sampled using temperatures ranging between 0.0 and 1.0, and their corresponding win rates are illustrated in Figure~\ref{fig:frontier-tldr-main} (right). The DPO, PPO, and Preferred-FT methods are all built upon fine-tuning the identical GPT-J SFT model\footnote{\url{https://huggingface.co/CarperAI/openai_summarize_tldr_sft}}. Our results indicate that DPO achieves a win rate of about 61\% at a temperature setting of 0.0, outperforming PPO, which attains around 57\% at its optimal temperature of 0.0. Furthermore, DPO secures a greater maximum win rate compared to the optimal $N$-sample baseline. It is important to mention that DPO's $\beta$ hyperparameter was not extensively tuned in these experiments, suggesting that the observed performance could be a conservative estimate of DPO's capabilities. In addition, DPO shows considerably more stability across different sampling temperatures than PPO, whose results tend to regress towards the base GPT-J model as temperature increases. The Preferred-FT method shows minimal improvements over the initial SFT baseline. Section~\ref{sec:human-judgments} presents a direct comparison between DPO and PPO via human evaluation, where DPO samples generated at temperature 0.25 were favored 58\% of the time relative to PPO samples at the same temperature.

For the single-turn dialogue experiments, we assess each approach using the subset of the Anthropic HH dataset \citep{bai2022training} test split that includes a single round of human-assistant interaction. For GPT-4-based evaluations, we measure win rates by referencing the preferred completions from the test set across different methods. As a standard SFT model is unavailable for this particular task, our pipeline begins with a pre-trained Pythia-2.8B model. We then apply Preferred-FT on the selected completions to create a reference model whose outputs remain within the model’s distribution, following which we train using DPO. For comparison, we consider the best completion among 128 Preferred-FT generations (the performance of the Best of $N$ method saturates at $N=128$ for this evaluation, see Appendix Figure~\ref{fig:best-of-n}) as well as a two-shot prompted baseline using Pythia-2.8B. The results indicate that DPO matches or exceeds the highest-performing temperature settings for each strategy. Additionally, we include results for an RLHF model trained with PPO on the Anthropic HH dataset\footnote{\url{https://huggingface.co/reciprocate/ppo_hh_pythia-6B}} from a recognized implementation\footnote{\url{https://github.com/CarperAI/trlx/tree/main/examples/hh}}, but despite extensive tuning of prompts and sampling temperatures, this RLHF model does not surpass the baseline performance of Pythia-2.8B. Drawing from our findings in TL;DR and acknowledging that both Best of 128 and PPO target the same reward function, we treat Best of 128 as an approximate indicator of PPO-level outcomes. In summary, DPO emerges as the only computationally practical method that achieves gains over the preferred completions on the Anthropic HH dataset and offers performance that is on par with or better than the more resource-intensive Best of 128 approach. Lastly, as shown in Figure~\ref{fig:dialogue-main}, DPO attains peak performance in relatively few training iterations.

\subsection{Generalization to a new input distribution}

To further examine how PPO and DPO cope with distributional changes, we assess the performance of the policies derived from our Reddit TL;DR summarization study on a novel data distribution: news articles from the test split of the CNN/DailyMail dataset \citep{nallapati-etal-2016-abstractive}. We use the optimal sampling temperatures found on TL;DR (0 and 0.25). Table~\ref{tab:ood} presents these findings. The evaluation metric is the win rate of GPT-4, comparing model-generated summaries to ground-truth summaries in this dataset, utilizing the same GPT-4 (C) prompt as in Reddit TL;DR—except that "forum post" is substituted with "news article." On this distinct distribution, DPO maintains a considerable advantage over PPO. This experiment offers preliminary support for the conclusion that DPO policies can generalize comparably to PPO policies, despite DPO not utilizing the extra unlabeled Reddit TL;DR prompts that PPO employs.

\subsection{Validating GPT-4 judgments with human judgments}
\label{sec:human-judgments}
We implement a human subject study to assess the dependability of GPT-4's evaluative judgments, drawing from results of the TL;DR summarization experiments and testing two variants of GPT-4 prompting. The \textbf{GPT-4 (S)} (simple) prompt asks directly which summary better encapsulates the key information in the post. The \textbf{GPT-4 (C)} (concise) prompt incorporates a request for conciseness, as we observe GPT-4 has a bias toward lengthier, more repetitive summaries when presented with the \textbf{GPT-4 (S)} prompt. The full prompt texts are provided in Appendix~\ref{app:prompts}. Our study compares the top-performing (DPO, temp. 0.25), lowest-performing (PPO, temp. 1.0), and an intermediate-performing (SFT, temp. 0.25) model, each against PPO with its optimal greedy decoding temperature, thereby encompassing a range of summary qualities. The outcomes indicate that GPT-4's preferences under both prompts align with human preferences approximately as well as human raters agree with each other, making GPT-4 a suitable stand-in for human judgment in this context (due to a restricted number of human annotators, multiple human judgments were gathered only for the DPO and PPO-1 comparisons). Across all conditions, the \textbf{GPT-4 (C)} prompt typically yields win rates closer to those of human raters; consequently, this prompt is used for principal results in Section~\ref{sec:dpo-real-datasets}. Further methodological details, including the web-based rater interface and a roster of volunteer annotators, can be found in Appendix~\ref{app:human-study}.

\section{Discussion}
Preference-based learning constitutes an effective and scalable approach for developing capable and value-aligned language models. In this work, we have presented DPO, a streamlined methodology for training language models using preference data, eliminating the need for reinforcement learning. Instead of recasting preference learning as a reinforcement learning (RL) problem to leverage standard RL algorithms, DPO directly establishes a connection between language model policies and reward functions. This enables language models to be optimized to reflect human preferences using a straightforward cross-entropy objective, thus obviating both reinforcement learning and any compromise in generality. Remarkably, DPO matches or surpasses the performance of established RLHF techniques—such as those based on PPO—with minimal hyperparameter tuning required, thereby substantially lowering the barrier to training additional models from preference data.

\textbf{Limitations \& Future Work.} The findings presented prompt a range of directions for subsequent research. One pertinent question is the extent to which the DPO-trained policy extrapolates beyond its training distribution, particularly in relation to models trained with explicit reward signals; our preliminary experiments indicate comparable generalization to PPO-based policies, but systematic investigations remain necessary. It is also of interest to assess whether DPO’s self-labeling training mechanism can exploit unlabeled prompts as effectively as other methods. Another avenue concerns understanding the dynamics of reward over-optimization in the DPO context, specifically whether the modest performance drop observed in Figure~\ref{fig:dialogue-main}-right reflects such effects. Additionally, although current evaluations focus on models with up to 6B parameters, expanding DPO to models with significantly larger scales constitutes a promising line of inquiry. Regarding evaluation protocols, our experiments reveal that win rates derived from GPT-4 can be sensitive to prompt design; optimizing the use of automated evaluators to ensure reliable assessments warrants future examination. Lastly, DPO holds promise for diverse generative modeling tasks beyond language modeling from preference data, opening further potential applications.